\documentclass[a4paper,10pt]{report}\input{../0.Préambule/Préambule_cours_prof}\begin{document}

\newpage
\section*{Activité : Pourcentage et conversion}
\addcontentsline{toc}{section}{Activité : Pourcentage et conversion}

\textbf{Consigne : }Vous devez vous mettre par deux et faire les exercices suivants sur une copie double que vous remettrez à la fin de l'heure. Sur la copie double je veux les deux écritures (une en couleur bleue, l'autre en couleur noire).

\begin{exer}

\begin{enumerate}
\item Lors d'une élection, dans une commune où 480 votes ont été exprimés, une candidate a obtenu $11,25\%$ des voix. Calcule le nombre de personnes qui ont voté pour elle.
\item Pour la même élection, un autre candidat a obtenu $132$ voix. Calcule le pourcentage de votes exprimés pour ce candidat.
\end{enumerate}
\end{exer}

\begin{exer}
Sur 720 000 jeunes sortis du système éducatif $8\%$ sont sans diplôme et $13\%$ ont au moins réussi le bac.

\medskip
\noindent Quatre ans plus tard, $32\%$ des \og sans diplôme \fg{} et $18\%$ des simples bacheliers sont au chômage.

\medskip
\noindent Calcule le nombre de chômeurs de chaque catégorie.
\end{exer}

\begin{exer}
On compte environ $\np{25 823 000}$ actifs en France en $2016$.
\begin{enumerate}
\item Sachant qu'il y  a environ $2,8\%$ d'agriculteurs, quel est leur nombre approximatif ?
\item Le nombre de personnes travaillant en $2016$ dans la construction est d'environ $\np{1 704 300}$ personnes. Calcule leur pourcentage à $0,1\%$ près par rapport au nombre d'actifs.
\end{enumerate}
\end{exer}

\begin{exer}
Introduit en Australie en $1935$ pour lutter contre les insectes rongeant la canne à sucre, le crapaud buffle, qui est venimeux, ravage désormais la faune locale.
\begin{enumerate}
\item La taille des $100$ spécimens introduits à l'origine était au maximum de $14$cm mais un spécimen de $38$cm a été capturé en $2007$. De quel pourcentage sa taille a-t-elle augmenté ?
\item Une estimation actuelle donne une population de crapauds buffles en Australie de l'ordre de $200$ millions d'individus. De quel pourcentage leur nombre a-t-il augmenté par rapport à $1935$ ?
\end{enumerate}
\end{exer}

\begin{exer}

\medskip
\begin{enumerate}
\item Convertir les vitesses ci-dessous en km/h.
\begin{center}
\renewcommand{\arraystretch}{1.5}
\begin{tabularx}{\linewidth}{|*{3}{>{\centering \arraybackslash}X|}}
\hline
\rowcolor{gray!20} Objet / Personne / Animal & Vitesse & Vitesse (en km/h)\\\hline
Randonneur expérimenté & 3s pour 5 m &\\\hline
Guépard en chasse & 270 m en 9s &\\\hline
Faucon pèlerin en piqué & 150 m en 3 s &\\\hline
Escargot motivé & 12 cm par minute &\\\hline
Boeing 747 au décollage & 69 m par seconde &\\\hline
Record d'Usain Bolt en 2009 & 9,58 s sur 100 m &\\\hline
Espadon Voilier & 11,2 km en 6 min &\\\hline
Boa Constrictor & 23 m en 4 min &\\\hline
Limace sportive & 30 min par mètre &\\\hline
Lièvre en pleine course & 1,2 km/min &\\\hline
Lévrier en cynodrome & 20,6 m/s &\\\hline
Autruche en fuite & 16,2 km en 10 min &\\\hline
Record du TGV en 2007 & 638 m en 4s &\\\hline
Vitesse d'une tornade & 130 m par seconde &\\\hline
Rotation de la terre à l'équateur & 464 m/s &\\\hline
Rotation de la Terre autour du Soleil & 29,8 km/s &\\\hline
\end{tabularx}
\end{center}	
\item Découper les vignettes associées en les rangeant dans l'ordre croissant.
\end{enumerate}

\noindent \includegraphics[scale=1]{../40.Image/Act-Pourcentage_et_conversion_1.png}
\end{exer}

\end{document}